\section{Research Objectives}

The desired outcome of this thesis is to improve upon \gls{UHF} methodology with steady-steady acquisitions for two neuroscience application namely \gls{fMRI} (introduced in \cref{sec:fmri} ) and \gls{DMI} (introduced in\cref{sec:DMI}).
As first part of the thesis, the acquisition speed of the recently proposed passband-\gls{bSSFP} \gls{fMRI} has to be improved to the point of the popular \gls{GE}-\gls{EPI}. 
The \gls{bSSFP} with its short readout and spin-echo behavior aims to provide an alternative for high resolution \gls{fmri} which suffer from the 

the gain in nominal spatial resolution is lost with macrovascular specificity and distortion at \gls{UHF}. 

The current \gls{GE}-\gls{EPI} is quite mature and optimized with many acceleration techniques which exploit receiver sentivity encoding, RF encoding in term of simultanius-multi-slice excitationm, low spatial frequency of phase with partial fourier.

The reconstruction of 3D undersampled non-Cartesian data is notoriously diffcult inverse problem

3D encoding is desired to reduce the inflow effects



The challenges at \gls{UHF} such as the highly off-resonant fat signal has to be suppressed for the low bandwidth spiral readout and the residual spatial off-resonance induced blurring has to be adequate corrected using deblurring algorithm. As several hundreds of volumes are acquired for the functional studies, the computation cost of the reconstruction model has to be reduced.

with the efficient \gls{NUFFT} algorithm for transforming non-uniform grid, the segmented $B_0$ effects for deblurring and proper pre-conditioning to speed up the convergence

In previous \gls{bSSFP} functional studies, highly segmented EPI17,18 and spiral readouts16 have been proposed to improve the speed of bSSFP acquisition. In publication 1, we aim to further improve speed and image quality by combining  parallel imaging and correction techniques for system imperfections.





 The favorable $T_2/T_1$ of deuterated metabolites yielded a much higher SNR improvement with \gls{bSSFP} compared to spoiled \gls{GRE} sequences
 
 Due to the significant spatial $B_0$ inhomogeneity in the human head at 9.4 T, phase-cycling was employed to reduce the off-resonance sensitivity, albeit with a moderate reduction in \gls{SNR}. 
 
 
 Glucose-DMI investigations in the human brain have reported spatio-temporal resolutions ranging from 3 mL in 10 min at 9.4 T to 14 mL in 39 min at 3 T2,9–11
 
 
 SNR optimality of the reference protocol is evident from the higher spatio-temporal resolution achieved in this study compared to the previous \gls{DMI}. 